\documentclass[11pt,a4paper]{article}
\usepackage[margin=19mm,headheight=15pt]{geometry}
\usepackage{fontspec}
\setmainfont{Noto Serif}
\usepackage{amsmath,amssymb,mathtools}
\usepackage{xcolor}
\usepackage{fancyhdr}
\usepackage{enumitem}
\usepackage{microtype}
\setlength{\parindent}{0pt}
\setlength{\parskip}{6pt}
\definecolor{navy}{HTML}{183153}
\definecolor{lightgray}{HTML}{F3F5F7}
\pagestyle{fancy}\fancyhf{}
\lhead{\small ME 5309 -- Mid-Semester Examination}
\rhead{\small ME 5309}
\cfoot{\thepage}
\begin{document}
\begin{center}
{\Large\bfseries ME 5309 -- Mid-Semester Examination}\\[3pt]
{\large Answer Booklet}\\[8pt]
\end{center}
\textbf{Instructions:} Attempt any three questions. Each question carries 10 marks. Show working where appropriate.
\section*{Question 1 \hfill [10 marks]}
\subsection*{Q1(a) \hfill [4 marks]}

(i) Pixel pitch p = 5 microm = 0.005 mm.
Hence:\[\begin{aligned}f_x = f/p = 16/0.005 = 3200 pixels\end{aligned}\]
(ii) Sensor width = 1280(0.005) = 6.40 mm.
At standoff Z = 400 mm, belt width W = Z(6.40)/16 = 160 mm.
\[\begin{aligned}Scale = 160/1280 = 0.125 mm/pixel\end{aligned}\]\par\smallskip\textit{Working:} $5\,\mu\mathrm{m}=5\times10^{-3}\,\mathrm{mm}=0.005\,\mathrm{mm}$. So $16/0.005=16\times200=3200$ pixels. For the sensor, $1280\times0.005=6.4$ mm. Then $W=(400/16)\times6.4=25\times6.4=160$ mm, and $160/1280=0.125$ mm/pixel.

\subsection*{Q1(b) \hfill [3 marks]}

ADC step = 5/2\textasciicircum{}12 = 0.0012207 V = 1.2207 mV.
Sensor sensitivity is 10 mV/N, so ΔF = 1.2207/10 = 0.12207 N.
The smallest resolvable force change is approximately 0.122 N.\par\smallskip\textit{Calculation:} $2^{12}=4096$ levels. One step is $5/4096=0.0012207$ V, i.e. $1.2207$ mV. Dividing by the sensitivity, $1.2207\,\mathrm{mV}/(10\,\mathrm{mV/N})=0.12207$ N. The mV units cancel, leaving N.

\subsection*{Q1(c) \hfill [3 marks]}

For zero-order hold, F = e\textasciicircum{}(AT) and G = integral\_0\textasciicircum{}T e\textasciicircum{}(Aτ)B dτ.
Here:\[\begin{aligned}A = -2, B = 1, T = 0.1 s\end{aligned}\]
Thus:\[\begin{aligned}F = e^(-0.2) = 0.818731\\ G = (1-e^(-0.2))/2 = 0.0906346\end{aligned}\]
Therefore:\[\begin{aligned}x_{k+1} = 0.818731 x_k + 0.0906346 u_k\end{aligned}\]\par\smallskip With $AT=(-2)(0.1)=-0.2$, $F=e^{-0.2}=0.8187308$. For $G$, $(1-F)/2=(1-0.8187308)/2=0.1812692/2=0.0906346$. Thus the input contribution is positive.

\section*{Question 3 \hfill [10 marks]}
\subsection*{Q3(a) \hfill [2 marks]}

Finite differences are poor mainly because the encoder measures angle in degrees instead of radians, and because the velocity signal is always delayed by one full second.

\subsection*{Q3(b) \hfill [2 marks]}

The Tustin coefficients are y[k]=0.881765y[k-1]+0.0591175(x[k]+x[k-1]).
The DC gain is 0.0591175/(1-0.881765)=0.5.

\subsection*{Q3(c) \hfill [4 marks]}

A-LC=[[-ℓ₁,1], [-ℓ₂,-2]].
The desired polynomial is (s+10)² = s²+20s+100.

\subsection*{Q3(d) \hfill [2 marks]}

The observer is simpler because it only smooths the measured velocity.
The filter is safer when the model is poor.

\section*{Question 4 \hfill [10 marks]}
\subsection*{Q4(a) \hfill [2 marks]}

Closed-loop characteristic polynomial is s²+(a+bk₂)s+(c+bk₁)=s²+4s+8, giving poles -2±2j.
At DC, θ/r = bNff/(c+bk₁), so unity gain requires Nff=(c+bk₁)/b=(6+2)/1=8.\par\smallskip Comparing coefficients with $s^2+4s+8$ gives $a+bk_2=4$ and $c+bk_1=8$. The feedforward gain is chosen from the nominal plant so the reference scaling is unity; here this gives $N_{ff}=8$. The closed-loop roots of $s^2+4s+8$ are $-2\pm2j$.

\subsection*{Q4(b) \hfill [2 marks]}

\[\begin{aligned}\theta_{ss} = 0.75\times8/(6+2)=0.75,\\ \text{the percentage tracking error is 25\%}\end{aligned}\]

\subsection*{Q4(c) \hfill [4 marks]}

With ξ̇=r-θ and u=-k₁θ-k₂θ̇+K\_Iξ, the homogeneous closed-loop characteristic polynomial is s³+(a+bk₂)s²+(c+bk₁)s+bK\_I.
Desired polynomial: (s+4)[(s+2)²+4]=(s+4)(s²+4s+8)=s³+8s²+24s+32.
For a=2,b=1,c=6, matching coefficients gives k₂=6, k₁=18 and K\_I=32.\par\smallskip First multiply $(s+4)(s^2+4s+8)=s^3+8s^2+24s+32$. Comparing this with the augmented characteristic polynomial gives $k_2=6$, $k_1=18$, and $K_I=32$. I checked the coefficients in descending powers of $s$.

\subsection*{Q4(d) \hfill [2 marks]}

\vspace{2mm}

\end{document}